\subsection{Finite atomic algebras and a diffuse limit}
\label{p:sec:structure}

The model also distinguishes the spectral information in finitely many
polar conjugates from that in the full orbit. This distinction follows
from the explicit weight and the rotation, independently of the
convolution argument.

\begin{proposition}\label{p:prop:joint-structure}
For the operator $T$ constructed above, its polar decomposition is
$T=U|T|$, with $|T|=\pi(f)$ and $\ker|T|=\{0\}$.
The spectral measure of $|T|$ is countably atomic. For every finite
$F\subset\Z$, the abelian von Neumann algebra
\[
 \mathcal A_F=W^*(U^n|T|U^{-n}:n\in F)
\]
is atomic. The full algebra
\[
 \mathcal A=W^*(U^n|T|U^{-n}:n\in\Z)
\]
is diffuse.
\end{proposition}

\begin{proof}
Since $f>0$ almost everywhere, $T^*T=\pi(f)^2$ has zero kernel, and
$|T|=\pi(f)$. The unitary $U$ is therefore its polar factor.
On the conull set where the sum in \eqref{p:eq:weights} has finite
support, the value of $f$ is determined by a finite subset of the
countable coordinate set. There are only countably many such subsets.
Consequently $f$ takes values in a countable subset of $(0,1]$, up to a
null set. Its nonnull level sets give countably many spectral projections
whose sum is $1$, which proves the assertion about spectral type.

The representation convention gives
\[
 U\pi(g)U^*=\pi(g\circ\sigma^{-1}).
\]
For a finite $F$, the tuple $(f\circ\sigma^{-n})_{n\in F}$ has
countable range on a common conull set. The nonnull joint level sets
therefore give an atomic partition of the generated algebra
$\mathcal A_F$. This includes $F=\varnothing$, when the generated
unital algebra is scalar.

For the infinite family, $\mathcal A$ is an invariant abelian subalgebra
of $\pi(L^\infty(X))$. The rotation is totally ergodic. Indeed, for any
integer $q\ne0$ and any nonconstant finite-support character indexed by
$v=(v_i)$, rational independence implies
$q\sum_i v_i\theta_i\notin\Z$. The Fourier argument used in
Section~\ref{p:sec:scales} thus proves ergodicity of $\sigma^q$.

Suppose $e$ were an atom of $\mathcal A$. Its integer translates under
$\operatorname{Ad}U$ are atoms of equal positive trace. Distinct atoms
in an abelian algebra are orthogonal, so only finitely many can occur.
The orbit of $e$ is therefore finite, and for some positive integer $q$,
$U^q eU^{-q}=e$. Because $e\in\pi(L^\infty(X))$, total ergodicity
forces $e=1$. An abelian algebra whose identity is an atom is scalar.
But $\mathcal A$ contains the nonconstant multiplication operator $\pi(f)$: a positive
almost-everywhere constant would have finite logarithmic integral,
contrary to Lemma~\ref{p:lem:quasinilpotent}. Hence $\mathcal A$ has
no atoms, as claimed.
\end{proof}

The factor construction, the norm and determinant calculation, the
exclusion of invariant projections, and
Proposition~\ref{p:prop:joint-structure} prove
Theorem~\ref{p:thm:model}, with $c_i=C_j$ for $i\in J_j$.
The diffuseness statement is only about $\mathcal A$; it does not
identify $\mathcal A$ with all of $\pi(L^\infty(X))$.
